\section*{अध्याय \ref{ch:g9:alternating-current} --- प्रत्यावर्ती धारा}

\begin{solution}{exo:g9:alternating-current:1}
DC टिके हुए \omterm{def:g8:voltage-voltmeter:voltage}{विभवांतर} के नीचे हमेशा एक ही ओर बहती है --- यानी \omterm{def:g3:first-electric-circuit:battery}{बैटरी} का
राज। और AC झूलते हुए \omterm{def:g8:voltage-voltmeter:voltage}{विभवांतर} के नीचे समय-समय पर उलटती रहती है ---
यानी सॉकेट का राज।
\end{solution}

\begin{solution}{exo:g9:alternating-current:2}
$T$: एक पूरे चक्र की \omterm{def:g2:measuring-time:duration}{अवधि}; $f$: हर सेकंड के चक्र; और $f = 1/T$ तथा
$T = 1/f$. \qty{50}{Hz} के लिए: $T = \qty{0.02}{s}$. और
$T = \qty{0.01}{s}$ के लिए: $f = \qty{100}{Hz}$.
\end{solution}

\begin{solution}{exo:g9:alternating-current:3}
सौ बार: पचास में से हर चक्र पर दो उलटाव।
\end{solution}

\begin{solution}{exo:g9:alternating-current:4}
$T = 4 \times 5 = \qty{20}{ms} = \qty{0.02}{s}$; $f = 1/0.02 =
\qty{50}{Hz}$.
\end{solution}

\begin{solution}{exo:g9:alternating-current:5}
\qty{230}{V} \omterm{prop:g9:alternating-current:effective}{प्रभावी विभवांतर} है --- वह टिका हुआ DC मान जो उतना
ही गरम करता, और जिसे AC \omterm{def:g8:voltage-voltmeter:voltmeter}{वोल्टमीटर} पढ़ते हैं; जबकि \qty{325}{V} शिखर है,
यानी वह चोटी जिसे वह झूला सचमुच छूता है, दोनों ओर, हर सेकंड पचास बार।
\end{solution}

\begin{solution}{exo:g9:alternating-current:6}
गरम होने को दिशा की कोई परवाह नहीं --- कुंडली किसी भी क़तार पर गरम हो
जाती है, इसलिए केतलियाँ वह झूला सीधे पी लेती हैं। जबकि चार्ज होना उलटी
दिशा में चलता रसायन है, और रसायन एक टिकी हुई दिशा माँगता है: इसलिए ईंट
पहले उस झूले को अटलता में बदल देती है।
\end{solution}

\begin{solution}{exo:g9:alternating-current:7}
\omterm{def:g3:first-electric-circuit:battery}{बैटरी}: \qty{9}{V} पर एक सपाट रेखा। प्रत्यावर्ती आपूर्ति: $+9$ और
\qty{-9}{V} के बीच झूलती एक मुलायम तरंग, हर \qty{20}{ms} में एक चक्र।
एक नज़र वाली परख: सपाट बनाम लहरदार।
\end{solution}

\begin{solution}{exo:g9:alternating-current:8}
घूमती मशीनें क़ुदरतन AC बनाती हैं; और ट्रांसफ़ॉर्मर को सिर्फ़ AC ही
खिलाती है, जिसके बिना बिजली ऊँचे \omterm{def:g8:voltage-voltmeter:voltage}{विभवांतर} पर देश पार नहीं कर सकती थी
और घरों में सधे हुए \omterm{def:g8:voltage-voltmeter:voltage}{विभवांतर} पर नहीं घुस सकती थी।
\end{solution}

\begin{solution}{exo:g9:alternating-current:9}
धीमे पैडल पर डाइनेमो के हर सेकंड के चंद चक्र लैंप को साफ़-साफ़ स्पंदित
कर देते हैं --- हर शिखर एक कौंध, हर शून्य एक गड्ढा; और तेज़ पैडल उन
स्पंदों को इतना गुणा कर देता है कि आँख उन्हें आपस में मिला देती है।
सिनेमा की $24$ तसवीरें आँख के मिलाने की मोटी दर तय कर देती हैं: क़रीब
\qty{25}{Hz} के प्रकाश-स्पंदों से नीचे टिमटिमाहट खलती है; और ऊपर दमक
टिकी हुई पढ़ी जाती है।
\end{solution}

\begin{solution}{exo:g9:alternating-current:10}
शिखर क़रीब \qty{325}{V} पर --- और हर सेकंड उनमें से सौ, यानी पचास धन
और पचास ऋण।
\end{solution}

\begin{solution}{exo:g9:alternating-current:11}
$T = 1/60 \approx \qty{0.017}{s}$; और शिखर $\approx 120 \times 1.4
\approx \qty{170}{V}$. हीटर और लैंप बेपरवाही से ढल जाते हैं; आधुनिक
चार्जर अपने लेबलों पर लिखा ``100--240 V, 50/60 Hz का परास'' पढ़ लेते
हैं और ख़ुशी से
बदल लेते हैं --- एतराज़ करने वाले तो पुराने, एक ही \omterm{def:g8:voltage-voltmeter:voltage}{विभवांतर} वाले उपकरण
हैं, और साथ में वह हर \omterm{ex:g6:circuits-and-safety:motor}{मोटर} वाली घड़ी जो जाल के चक्र गिनती है: वह साठ पर
तेज़ चलने लगती।
\end{solution}

\begin{solution}{exo:g9:alternating-current:12}
दोनों शिखरों पर चमकने से हर सेकंड $2 \times 50 = 100$ स्पंद मिलते हैं।
और $25$ फ़्रेम वाला कैमरा हर \qty{0.04}{s} पर नमूना लेता है --- यानी
हर फ़्रेम पर चार स्पंद, पर लुढ़कते हुए खुलने के हर बिंदु पर एक-सा
क़तारबद्ध नहीं: कैमरे की लय और लैंप की लय के बीच की वह ज़रा-सी बेमेली
ऐसी पट्टियाँ पोत देती है जो एक के बाद एक फ़्रेमों पर सरकती जाती हैं ---
यानी दो घड़ियों के बीच की टकराहट। और \qty{60}{Hz} पर ($120$ स्पंद) $25$
फ़्रेमों के सामने वह बेमेली बदल जाती है, इसलिए पट्टियों की चौड़ाई और
\omterm{def:g7:motion-average-speed:formula}{चाल} भी बदल जाती है --- वह नमूना स्थानीय जाल का भेद खोल देता है।
\end{solution}

\begin{solution}{pb:g9:alternating-current:1}
\textbf{1.} $2.25 \times 2 = \qty{4.5}{V}$ पर DC: यानी जेब वाली चपटी
\omterm{def:g3:first-electric-circuit:battery}{बैटरी}।
\textbf{2.} AC; $T = \qty{20}{ms}$, $f = \qty{50}{Hz}$; और शिखर
$3 \times 2 = \qty{6}{V}$.
\textbf{3.} \qty{4.5}{V} की DC, \emph{उलटी} जुड़ी हुई --- तार
आपस में बदले हुए: \omterm{met:g9:alternating-current:scope}{दोलनदर्शी} पर बेज़रर, और बहीखाते के लिए सिखाने वाली।
\textbf{4.} AC; $T = \qty{10}{ms}$, $f = \qty{100}{Hz}$; और शिखर
\qty{3}{V}। उसके हर सेकंड सौ दोहरे स्पंद B की पचास-चक्र वाली टिमटिमाहट
से कहीं आराम से आँख के मिलाने की दर के पार बैठते हैं।
\textbf{5.} \omterm{def:g8:voltage-voltmeter:voltmeter}{वोल्टमीटर} प्रभावी मान बताता है --- यानी गरम करने के लिहाज़
से DC वाला बराबरी का मान --- जो ज्यावक्र के लिए क़रीब शिखर बटा $1.4$
पर बैठता है: $6 \div 1.4 \approx \qty{4.3}{V}$. दोनों अंक एक ही झूले
का बयान करते हैं।
\textbf{6.} सिर्फ़ आपूर्ति A (सही ओर से जुड़ी हुई --- और दोबारा जोड़ने
के बाद C भी)। AC वाली आपूर्तियों तथा किसी भी \omterm{def:g3:first-electric-circuit:battery}{बैटरी} के बीच एक बदलने
वाला यंत्र चाहिए --- चार्जर की दिष्टकारी ईंट।
\textbf{7.} मेन के हर सेकंड के पचास झूले दिल की अपनी विद्युत लय के पास
पड़ते हैं और उसे बिगाड़ सकते हैं: \omterm{def:g8:intensity-ammeter:intensity}{ऐम्पियर} दर \omterm{def:g8:intensity-ammeter:intensity}{ऐम्पियर} वह झूला टिकी हुई
क़तार से ज़्यादा ख़तरा है।
\textbf{8.} ज्यावक्र --- यानी वह पूर्ण झूला, जिसे किसी चुंबकीय क्षेत्र
में अटल \omterm{def:g7:motion-average-speed:formula}{चाल} से घूमती कुंडली क़ुदरतन खींच देती है: अगले अध्याय का
प्रत्यावर्तित्र।
\textbf{9.} हर $4$ खानों में एक चक्र (\qty{20}{ms}); और शिखर ऊपर तथा
नीचे $325 \div 100 = 3.25$ खाने।
\textbf{10.} $f = 1/0.02 = \qty{50}{Hz}$; और प्रभावी $\approx 325
\div 1.4 \approx \qty{230}{V}$ --- वही लेबल, पर्दे से बरामद।
\textbf{11.} \qty{230}{V} पर एक सपाट रेखा --- यानी बराबर ऊष्मा-शक्ति
वाली अटलता।
\textbf{12.} जैसे: ``सपाट रेखा \omterm{def:g3:first-electric-circuit:battery}{बैटरी} की अटलता है --- एक दिशा, एक मान।
तरंग जाल का झूला है --- दिशा और मान, दोनों अनंत चक्र में। और दोनों के
बीच इंसाफ़ वाली विनिमय दर \omterm{prop:g9:alternating-current:effective}{प्रभावी विभवांतर} है: बराबर \omterm{def:g5:heat-and-insulation:heat}{ऊष्मा}, बराबर
मोल।''
\end{solution}
